\section{Proof of the counterexample}
\label{sec:assembly}

We now combine the geometric Morse comparison, the two manifold constructions, and the perturbation of the involution.

\begin{proof}[Proof of Theorem~\ref{thm:main}]
Fix an integer $N\geq704$ and take the simply connected Kähler manifold $X$ of Section~\ref{sec:geometry}. Let $Y$ be the Kähler surface of Proposition~\ref{prop:surface}, and put
\[
 M=X\times Y,\qquad C=X\times\CP^1,
\]
with the product symplectic form on $M$. Both $M$ and $C$ are simply connected. Their real dimensions are $2N+4$ and $2N+2$, respectively, so Theorem~\ref{thm:smale} gives
\[
 \SM(M)=\lambda(M),\qquad \Morse(C)=\lambda(C).
\]
Choose a Morse function on $C$ with exactly $\lambda(C)$ critical points. Corollary~\ref{cor:product-realization}, applied to the four fixed components of $\id_X\times g_Y$, provides a smooth one-periodic Hamiltonian $H$ whose time-one map has exactly these $4\lambda(C)$ fixed points, all nondegenerate. Every corresponding periodic orbit is contractible because $M$ is simply connected. Proposition~\ref{prop:gap} now gives
\[
 \#\Fix_0(\phi_H^1;H)=4\lambda(C)=\lambda(M)-16=\SM(M)-16.
\]
More explicitly, the construction has
\[
 \SM(M)=456N-2056,\qquad
 \#\Fix(\phi_H^1)=456N-2072.
\]
Taking $N=704$ gives the dimension and counts stated in the theorem.
\end{proof}

\begin{remark}[Comparison with the integral Floer bound]\label{rem:periodic-bound}
The distinction between integer and cyclic gradings can be seen numerically in this family. The \emph{minimal Chern number} $\nu$ is the positive generator of the subgroup $\langle c_1(TM),\pi_2(M)\rangle\subset\Z$ when that subgroup is nonzero. An exceptional line in $Y$, regarded as a sphere in $M$ with its $X$-coordinate fixed, has Chern number $2-1=1$: its tangent line has degree $2$ and its normal line has degree $-1$. Hence $\nu=1$ for $M$.

The integral bound of Bai--Xu \cite[Theorem~A]{BaiXu2022}, equivalently \cite[Corollary~1.2]{BaiXu2026}, therefore uses the two groups
\[
 \overline H_{\bar i}(M;\Z)=\bigoplus_{k\equiv i\ (\mathrm{mod}\ 2)}H_k(M;\Z),
 \qquad \bar i\in\Z/2.
\]
In each parity the torsion is
\[
 (\Z/2)^{16(N-3)}\oplus(\Z/3)^{32(N-5)},
\]
by summing the profiles of Section~\ref{sec:counts}. Its minimum number of generators is $32(N-5)$. Thus their lower bound evaluates to
\[
 b(M)+2\sum_{\bar i\in\Z/2}d\bigl(\Tor\overline H_{\bar i}(M;\Z)\bigr)
 =8(41N-183)+128(N-5)=456N-2104,
\]
where $d$ denotes the minimum number of generators. Our Hamiltonian has $456N-2072$ fixed points, thirty-two more than this bound. The counterexample therefore distinguishes the integer-graded geometric Morse count from the established cyclically graded integral count.
\end{remark}
